\documentclass[11pt]{article}
\usepackage[a4paper,margin=25mm]{geometry}
\usepackage{amsmath,amssymb,amsthm,booktabs,microtype}
\usepackage[T1]{fontenc}
\usepackage{lmodern}
\usepackage[dvipsnames]{xcolor}
\usepackage[colorlinks=true,linkcolor=MidnightBlue,citecolor=MidnightBlue,urlcolor=MidnightBlue]{hyperref}
\usepackage{fancyhdr}
\pagestyle{fancy}\fancyhf{}
\fancyhead[L]{\small The multi-marginal converse of Minty's theorem}
\fancyhead[R]{\small Dev D. Goyal}
\fancyfoot[C]{\thepage}
\setlength{\headheight}{14pt}
\hypersetup{pdftitle={The multi-marginal converse of Minty's theorem fails in the plane},
pdfauthor={Dev D. Goyal},
pdfsubject={Fully autonomous agentic research},
pdfkeywords={multi-marginal optimal transport, c-monotonicity, maximal monotone operators, Minty theorem, firmly nonexpansive mappings, Lean 4}}
\newcommand{\replaystatus}{A fresh replay on 31 October 2026 compiled all six project modules and the audit module against a Mathlib checkout built from source at the pinned revision. All 47 declarations defined in the project modules, including 18 generated or auxiliary declarations, depend only on \texttt{propext}, \texttt{Classical.choice} and \texttt{Quot.sound}. The exact rational checker also passed; the record is in \texttt{verification/}.}
\newtheorem{theorem}{Theorem}
\newtheorem{lemma}[theorem]{Lemma}
\newtheorem{corollary}[theorem]{Corollary}
\newcommand{\R}{\mathbb R}
\newcommand{\ip}[2]{\langle #1,#2\rangle}
\newcommand{\norm}[1]{\lVert #1\rVert}
\title{\textbf{The multi-marginal converse of Minty's theorem\\fails in the plane}}
\author{Dev D. Goyal}
\date{31 October 2026}
\begin{document}
\maketitle
\begin{abstract}
For the three-marginal quadratic cost on the real plane, we give three
integer-coordinate tuples that are $c$-monotone but admit
no compatible tuple whose coordinate sum is zero. Four necessary
inequalities have a positive weighted sum bounded below by the exact
constant $23/26$, as witnessed by a sum of four squares. Every maximal
$c$-monotone extension of these tuples therefore has a proper, nonconvex
sum image. This gives a negative answer to the implication asked about in
Remark~2.6 of Bartz, Bauschke, Phan and Wang, using their two-point
definition of $c$-monotonicity. The actual finite obstruction, cost-swap
equivalence, maximal-extension argument and nonconvexity conclusion are
formalized in Lean~4.
\end{abstract}

\paragraph{Verification status.} \replaystatus

\section{The question and the precise conclusion}

Let $H=\R^2$ with its ordinary Euclidean inner product, let $X=H^3$, and put
\[
 c_3(x_1,x_2,x_3)=\ip{x_1}{x_2}+\ip{x_1}{x_3}+\ip{x_2}{x_3},
 \qquad S(x)=x_1+x_2+x_3.
\]
For $K\subseteq\{1,2,3\}$, let $m_K(x,y)$ take coordinate $i$ from $x$ when
$i\in K$ and from $y$ otherwise. A set $\Gamma\subset X$ is
\emph{$c_3$-monotone} if, for every $x,y\in\Gamma$ and every such $K$,
\begin{equation}\label{eq:swap}
 c_3(m_K(x,y))+c_3(m_K(y,x))\le c_3(x)+c_3(y).
\end{equation}
This is two-point $c$-monotonicity. No claim about full cyclic
$c$-monotonicity is made. Maximality means maximality by inclusion among
all subsets of $X$ satisfying~\eqref{eq:swap}.

Writing $d_i=x_i-y_i$, expansion of the four costs gives
\begin{equation}\label{eq:cut}
 c_3(x)+c_3(y)-c_3(m_K(x,y))-c_3(m_K(y,x))
 =\left\langle\sum_{i\in K}d_i,\sum_{i\notin K}d_i\right\rangle.
\end{equation}
Consequently, for three marginals it is enough to check the three singleton
cuts:
\begin{equation}\label{eq:single}
 \ip{x_i-y_i}{S(x)-S(y)-(x_i-y_i)}\ge0,\qquad i=1,2,3.
\end{equation}
Complementary cuts give the same inner product. The empty and full cuts
give zero. This agrees with Definition~1.2 and Lemma~2.3 of~\cite{BBPW}.

For a $c_3$-monotone set $\Gamma$, Bartz et al. prove that
$S(\Gamma)=H$ implies maximal
$c$-monotonicity, together with several equivalent full-domain resolvent
conditions~\cite[Theorem~2.5]{BBPW}. Their Remark~2.6 asks whether
maximality implies this sum-surjectivity. The following theorem negates
that implication.

\begin{theorem}\label{thm:main}
There is a three-element $c_3$-monotone set $\Gamma_0\subset(\R^2)^3$ such
that every $c_3$-monotone set $\Gamma\supseteq\Gamma_0$ satisfies
\[
 0\notin S(\Gamma),\qquad S(\Gamma)\text{ is not convex}.
\]
There exists a maximally $c_3$-monotone such $\Gamma$. In particular,
maximal $c_3$-monotonicity does not imply $S(\Gamma)=\R^2$.
\end{theorem}

The maximal set in this theorem is obtained by Zorn's lemma. The explicit
object is the finite obstruction $\Gamma_0$; we do not assert that this
three-element set is itself maximal.

\section{The integer-coordinate obstruction}

Define $\Gamma_0=\{A,B,C\}$ by the following table.
\begin{center}
\begin{tabular}{c@{\quad}ccc@{\quad}c}
\toprule
Point & First marginal & Second marginal & Third marginal & Sum\\
\midrule
$A$ & $(-1,-2)$ & $(3,-12)$ & $(9,-2)$ & $(11,-16)$\\
$B$ & $(-2,-5)$ & $(0,-7)$ & $(-14,-2)$ & $(-16,-14)$\\
$C$ & $(2,0)$ & $(-4,15)$ & $(14,-1)$ & $(12,14)$\\
\bottomrule
\end{tabular}
\end{center}

The three values on the left of~\eqref{eq:single} are as follows.
\begin{center}
\begin{tabular}{c@{\qquad}rrr}
\toprule
Pair & $i=1$ & $i=2$ & $i=3$\\
\midrule
$A,B$ & $11$ & $57$ & $92$\\
$A,C$ & $50$ & $25$ & $9$\\
$B,C$ & $211$ & $4$ & $27$\\
\bottomrule
\end{tabular}
\end{center}
Every entry is strictly positive. Reversing a pair changes both factors
in the inner product by a sign, and equal pairs give zero. Thus the table
and~\eqref{eq:cut} prove $c_3$-monotonicity for every pair and every cut.

\begin{lemma}\label{lem:nozero}
No $U=(u,v,w)\in(\R^2)^3$ with $u+v+w=0$ is $c_3$-monotonically related
to all three points of $\Gamma_0$.
\end{lemma}
\begin{proof}
Write $u=(a,b)$, $v=(c,d)$ and $w=(e,f)$. Fix a seed point $P$ with sum
$s=S(P)$ and one marginal $i$. Since $S(U)=0$, the corresponding
inequality~\eqref{eq:single} is equivalent to
\begin{equation}\label{eq:qgeneral}
 \norm{U_i-P_i}^2+\ip{U_i-P_i}{s}\le0.
\end{equation}
Using $(P,i)=(A,1),(B,2),(B,3),(C,2)$ respectively, expansion gives the
four necessary inequalities
\begin{align*}
 q_{A1}&=a^2+b^2+13a-12b-16\le0,\\
 q_{B2}&=c^2+d^2-16c-49\le0,\\
 q_{B3}&=e^2+f^2+12e-10f-52\le0,\\
 q_{C2}&=c^2+d^2+20c-16d+79\le0.
\end{align*}
The zero-sum hypothesis gives $e=-a-c$ and $f=-b-d$. Direct polynomial
expansion then establishes the identity
\begin{align}\label{eq:certificate}
 4q_{A1}+3q_{B2}+4q_{B3}+8q_{C2}
 ={}&8\left(a+\frac c2+\frac14\right)^2
     +8\left(b+\frac d2-\frac12\right)^2\notag\\
   &+13\left(c+\frac{31}{13}\right)^2
     +13\left(d-\frac{42}{13}\right)^2
     +\frac{23}{26}.
\end{align}
The left side is nonpositive, since its four weights are positive.
The right side is at least $23/26>0$. This contradiction proves the lemma.
No numerical optimizer or numerical tolerance is used in this argument.
\end{proof}

\section{Maximality and a nonconvex sum image}

\begin{proof}[Proof of Theorem~\ref{thm:main}]
The preceding table proves that $\Gamma_0$ is $c_3$-monotone. If a
$c_3$-monotone superset $\Gamma$ contained a tuple of sum zero, that tuple
would be compatible with all of $\Gamma_0$, contradicting
Lemma~\ref{lem:nozero}. Thus $0\notin S(\Gamma)$ for every such superset.

Consider the family of all $c_3$-monotone supersets of $\Gamma_0$, ordered
by inclusion. It is nonempty. The union of a nonempty chain belongs to
the family: any two tuples in the union lie together in one chain member,
where every inequality~\eqref{eq:swap} holds. Zorn's lemma supplies a
maximal family member $\Gamma$. Any $c_3$-monotone proper superset of
$\Gamma$ would still contain $\Gamma_0$, contradicting this maximality.
Hence $\Gamma$ is maximal among all $c_3$-monotone sets.

Finally, the three sums satisfy the exact relation
\[
 28S(A)+173S(B)+205S(C)=0,
 \qquad 28+173+205=406.
\]
All three weights are positive. If $S(\Gamma)$ were convex, it would
contain their normalized convex combination zero. This contradicts the
already established exclusion. The argument applies to every
$c_3$-monotone superset, including the maximal one.
\end{proof}

\section{The simultaneous-extension obstruction}

The relation to Minty parametrization is explicit. If $x,y$ belong to a
$c_3$-monotone set and $S(x)=S(y)$, inequality~\eqref{eq:single} reads
$-\norm{x_i-y_i}^2\ge0$ for each $i$, hence $x=y$. Thus on
$D=S(\Gamma)$ there are well-defined component maps
$J_i(S(x))=x_i$. They satisfy
\[
 \sum_{i=1}^3J_i(s)=s,\qquad
 \norm{J_i(s)-J_i(t)}^2\le
 \ip{J_i(s)-J_i(t)}{s-t}.
\]
These are precisely the firmly nonexpansive inequalities. For the
three-point domain $D_0=\{S(A),S(B),S(C)\}$, the displayed data therefore
give a partition of the identity into three firmly nonexpansive maps.
There is no simultaneous extension of these maps to $D_0\cup\{0\}$ that
preserves both firm nonexpansiveness and the identity-sum condition.
Indeed, the three values at zero would form the tuple excluded by
Lemma~\ref{lem:nozero}. Equivalently, the maps $Q_i=2J_i-\operatorname{Id}$
are nonexpansive on $D_0$ and sum to $-\operatorname{Id}$, but cannot be
extended jointly to zero while retaining these properties.

This consequence concerns compatibility of the three extensions. It does
not assert a failure of the classical extension or Minty theorem for one
ordinary monotone operator.

\section{Prior work and historical limits}

The target is the implication in~\cite[Remark~2.6]{BBPW}, not the
different question about contact sets of mutually $c$-conjugate
potentials. Lim's multidimensional example~\cite[Proposition~1.10 and
Example~4.1]{Lim} has a hexagonal gap in its sum image, but Lim explicitly
shows in Step~5 that its contact set is not maximally $c$-monotone and can
be enlarged by the zero tuple. In the present construction, every
$c_3$-monotone extension is prevented from adding any zero-sum tuple.
The distinction is essential to answering the maximality question.

The 2025 survey~\cite{Survey} discusses the multimarginal setting and
cites~\cite{BBPW}; it does not supply the finite obstruction above.
Finite incompatibility certificates, constrained Lipschitz extension
obstructions, and sum-of-squares identities are established techniques.
For example, Dacorogna and Gangbo~\cite{DG} study finite obstructions for
bounded-strain extension problems, whose hypotheses differ from the
simultaneous firmly nonexpansive identity-partition condition here.
No novelty is claimed for those general mechanisms.
The simplex/difference-body extremizers of classical Jung geometry also
give the hexagonal radius ratio $2/3$; see
\cite[Lemma~2.9 and Corollary~2.10]{JungGeometry}. A separate
three-dimensional obstruction considered during this work is a short
lift of that known geometry. Its existence is not evidence of a new
extension method. Whether earlier constrained-section results already
yield the present maximality conclusion remains a historical question.

Searches in October 2026 of the arXiv, citing records and code
repositories found no earlier answer to the question of Remark~2.6;
Bartz listed a multi-marginal Minty characterization among open questions
in 2020 and 2021 conference talks. The search is bounded, no external
expert has assessed the result, and historical priority is not certified.

\section{Formal verification}\label{sec:lean}

The formal development uses Lean~4.19.0 and the following mathlib commit:
\begin{center}\small
\nolinkurl{c44e0c8ee63ca166450922a373c7409c5d26b00b}.
\end{center}
It defines the real coordinate space, the Euclidean pairing, the actual
cost-swap predicate and inclusion-maximality. It verifies the table,
the polynomial identity, exclusion of every real zero-sum tuple,
equivalence with all cost swaps, the chain-union and Zorn argument, and
the actual nonconvexity of the sum image. The principal endpoint is
\begin{center}\small
\nolinkurl{Minty.exists_maximal_cMonotone_nonconvex_sum}.
\end{center}
The module \texttt{Minty.AxiomAudit} enumerates every declaration whose
defining module is one of the six project modules, including private and
generated declarations, and rejects any axiom other than
\texttt{propext}, \texttt{Classical.choice} and \texttt{Quot.sound}. The
script \texttt{scripts/check\_axiom\_audit.py} checks that output, and
\texttt{lean/CORRESPONDENCE.md} maps each step above to its declarations.
The standard-library script \texttt{certificate/verify\_minty\_exact.py}
checks the table and the polynomial identity in exact arithmetic. The
interpretation through firmly nonexpansive maps in the previous section
is proved in the text and is not packaged separately in Lean.

\section{Limitations}\label{sec:limits}

The theorem concerns three marginals on the real plane and the two-point
notion of $c$-monotonicity of~\cite{BBPW}. It does not classify other
numbers of marginals or dimensions, assert minimality of the dimension,
or address full cyclic $c$-monotonicity or $c$-splitting potentials. The
maximal set is obtained by Zorn's lemma and is not described explicitly.
Lean verification establishes only the formal statements described in
Section~\ref{sec:lean}, and the literature comparison does not establish
historical priority.

\section*{Research process and AI disclosure}

This is fully autonomous agentic research. Dev D. Goyal is responsible
for publication. No expert validation or external peer review is
claimed, and the novelty assessment is provisional.

The question selection, the counterexample, the proof, the Lean
formalization, the literature searches and an earlier version of this
note were produced by autonomous AI research agents in an earlier
session. The records of that session state that it used an autonomous
OpenAI GPT-6 Astra workflow with delegated discovery and review roles,
and that one additional review used GPT-6.1 Sol. A numerical search
suggested the witness; no solver output enters the proof.

Opus 5.5 by Anthropic used with Claude Code assessed the package for this
release. It read Remark~2.6 and Definition~1.2 of~\cite{BBPW} and the
relevant parts of~\cite{Lim}, compared the Lean definitions of the cost,
the swap condition and maximality with them, rechecked the table, the
identity and the convex combination in exact arithmetic with its own
code, ran the fresh Lean replay, and revised this note. AI checks are not
human peer review.

\begin{thebibliography}{9}
\bibitem{BBPW}
S.~Bartz, H.~H.~Bauschke, H.~M.~Phan and X.~Wang,
\emph{Multi-marginal maximal monotonicity and convex analysis},
Mathematical Programming \textbf{185} (2021), 385--408;
published online in 2019.
\href{https://doi.org/10.1007/s10107-019-01433-9}{doi:10.1007/s10107-019-01433-9}.
Preprint \href{https://arxiv.org/abs/1901.03777v2}{arXiv:1901.03777v2}.

\bibitem{Lim}
T.~Lim, \emph{Maximal monotonicity and cyclic involutivity of
multiconjugate convex functions}, SIAM Journal on Optimization
\textbf{33} (2023), 2489--2511.
\href{https://doi.org/10.1137/23M1549201}{doi:10.1137/23M1549201}.
Preprint \href{https://arxiv.org/abs/2207.04830v4}{arXiv:2207.04830v4}.

\bibitem{Survey}
L.~De Pascale, A.~Kausamo and K.~Wyczesany,
\emph{60 years of cyclic monotonicity: a survey},
\href{https://arxiv.org/abs/2308.07682}{arXiv:2308.07682}.

\bibitem{DG}
B.~Dacorogna and W.~Gangbo, \emph{Extension theorems for vector valued
maps}, Journal de Mathematiques Pures et Appliquees (2006).
\href{https://doi.org/10.1016/j.matpur.2005.04.005}{doi:10.1016/j.matpur.2005.04.005}.

\bibitem{JungGeometry}
R.~Brandenberg and B.~Gonz\'alez Merino,
\emph{The asymmetry of complete and constant width bodies in general
normed spaces and the Jung constant}, 2015 preprint version.
\href{https://arxiv.org/abs/1412.8693v3}{arXiv:1412.8693v3}.
\end{thebibliography}
\end{document}
